\chapterimage{golden_gate.jpg}

\chapter{Capa 3: comunicación Internet(work)}\label{sec:layer3}

\begin{minipage}{0.4\linewidth}
\begin{center}
\begin{bytefield}{16}
\bitbox{16}{Layer 7: Application} \\
\bitbox{16}{Layer 4: Transport} \\
\bitbox{16}{\color{color1} Layer 3: Internet} \\
\bitbox{16}{Layer 2: Network (LAN)} \\
\bitbox{16}{Layer 1: Physical} \\
\end{bytefield}
\end{center}
\end{minipage}
\begin{minipage}{0.6\linewidth}
\begin{center}
\includegraphics[width=\linewidth]{network_layer3.pdf}
\end{center}
\end{minipage}

\vspace{-0.75cm}


% \section{Layer~3 overview}
\subsection*{Capacidades}

El \conceptRef{IP}{Internet Protocol (IP)} permite intercambiar datagramas 
(PDUs de la capa~3) entre dispositivos de distintas LANs (redes de área local).
% 
Su sistema de direccionamiento identifica los dispositivos de forma global y hace posible 
encaminar datagramas allí donde lo necesitemos. 

Los dispositivos que se comunican mediante IP pueden pertenecer a LANs
basadas en tecnologías diferentes y con distintos tipos de dirección MAC.
Esto tiene dos implicaciones:
\begin{itemize}
  \item Un datagrama que ``cabe'' en la MTU de la LAN de origen
    puede ser demasiado grande para la MTU de la LAN de destino (o de cualquier LAN intermedia). 
    IP define un sistema de fragmentación para dividir datagramas,
    que solo se reconstruyen en el destino final.\\[-0.25cm]
    
  \item Cuando enviamos un datagrama, conocemos la \conceptRef{IP}{dirección IP} del dispositivo de destino,
  pero no necesariamente su dirección MAC. Solo necesitaremos esa MAC si formamos parte de la LAN
  de destino. IP usa ARP para traducir direcciones IP conocidas
  a direcciones MAC del tipo de la LAN de destino.
\end{itemize}


Los datagramas pueden viajar por redes que no controlan 
ni el origen ni el destino, por lo que pueden perderse, desordenarse y duplicarse.
% 
La capa~3 no proporciona mecanismos para tratar esto; las capas superiores deben implementar
mecanismos de protección cuando sea necesario 
(\eg\ para implementar los flujos de datos de TCP).

\begin{center}
\includegraphics[width=0.9\linewidth]{protocols_layer3.pdf}
\end{center}


\vspace{-0.25cm}
\subsection*{Protocolos}

La versión~4 de \concept{IP} es \textit{el} protocolo usado en la capa~3, \ie,
es común a prácticamente todas las comunicaciones de Internet tal y como la conocemos hoy.
% 
Desde hace muchos años, Internet intenta hacer la transición hacia IPv6, 
pero el proceso no está completo y solo IPv4 proporciona cobertura global.

IPv4 delega en el protocolo ARP la traducción de direcciones IP conocidas 
a direcciones MAC dentro de cada LAN. ARP \textit{no} usa datagramas IP.

Los siguientes protocolos encapsulan sus PDUs en el payload (carga útil) de los 
datagramas IP:\\[-0.6cm]
\begin{itemize}
\item Transport Control Protocol (TCP) 
\item User Datagram Protocol (UDP)
\item Internet Control Message Protocol (ICMP)
\end{itemize}

\begin{remark}
Los comandos \inlineCode{ip address}, \inlineCode{ip route} y \inlineCode{ip neighbor} 
(entre otros) te permiten consultar varios aspectos de tu configuración IP.
\end{remark}



\section{Direccionamiento en la capa~3}\label{sec:layer3:addressing}

Las \conceptRef{IP address}{direcciones} IPv4 tienen $32$~bits de longitud. Lo más habitual es presentarlas a las personas en
formato \conceptRef{quad decimal}{quad decimal}, \eg\ 
\begin{center}
\otherBase{142.250.201.67}
\end{center}

Las direcciones IPv6 tienen $128$~bits ($12$~bytes) de longitud, \ie, 
son $4$~veces mayores que una dirección IPv4.
Se expresan en formato hexadecimal, \eg\ como sigue. Fíjate en cómo 
dos signos de dos puntos ``\otherBase{::}'' indican ``rellena con ceros'',
y ``\otherBase{::}'' solo puede aparecer una vez por dirección IPv6.
\begin{center}
\otherBase{fe80::f524:a73a:e946:7c02}
\end{center}

\begin{remark}
Cuando una dirección IPv6 se usa como parte de una URL, se rodea con corchetes,
\eg\ \url{https://[fe80::f524:a73a:e946:7c02]/}.
\end{remark}


\subsection{Direcciones públicas y reservadas}\label{sec:layer3:public_reserved}

La mayoría de las direcciones IP son \conceptRef{public/private IP}{públicas} y únicas en todo Internet,
\ie, la IP anterior tiene el mismo significado en todo el mundo: identifica un único dispositivo conectado.
% 
Muchas direcciones son \conceptRef{public/private IP}{reservadas/privadas} y solo pueden 
usarse dentro de un ámbito concreto (\eg\ una LAN) o en una situación concreta.
Algunas de las más comunes son las siguientes (hay 
\href{https://en.wikipedia.org/wiki/Reserved_IP_addresses#IPv4}{\underline{algunas más}}, 
también para IPv6):

\begin{center}
\begin{tabular}{ll}
\toprule
\textbf{Bloque de direcciones IP} & \textbf{Ámbito de la dirección} \\
\toprule
\otherBase{127.*.*.*} & Este ordenador (localhost) \\[0.1cm]
\otherBase{10.*.*.*} & Interno (privado) \\
\otherBase{172.16.*.*} & Interno (privado) \\
\otherBase{192.168.*.*} & Interno (privado) \\[0.1cm]
\otherBase{0.*.*.*} & Usos especiales \\
\otherBase{169.254.*.*}& Usos especiales \\
\otherBase{255.255.255.255} & Usos especiales\\
\bottomrule
\end{tabular}
\end{center}

\begin{exercise}\ \\[-0.5cm]
\begin{itemize}
\item ¿Cuántas direcciones IPv4 distintas hay (incluyendo las reservadas y privadas)? 
\item ¿Son suficientes ahora y en el futuro?
\item ¿Es problemático que una dirección como \otherBase{192.168.0.1} 
  la estén usando simultáneamente miles (probablemente millones) de ordenadores en Internet ahora mismo?
\item ¿Cuáles de estas corresponden a direcciones link local?
\end{itemize}
\end{exercise}

\subsection{Máscaras de red}\label{sec:layer3:netmasks}

Los $32$~bits de una dirección IP pueden dividirse en dos partes usando una \conceptRef{netmask}{máscara de red}:\\[-0.25cm]
% 
\begin{center}
\begin{bytefield}{32}
\bitheader{0-31}\\
\bitbox{10}{Network ID} & \bitbox{22}{Device ID} \\
\end{bytefield}
\end{center}
% 
La primera parte identifica la LAN o el grupo de LANs al que pertenece un dispositivo. 
La segunda parte identifica un único dispositivo dentro de la LAN. 
La máscara de red determina cuántos bits se asignan a cada parte. 

En la figura anterior, la máscara de red asigna 10 bits al Network ID y 22 al Device ID.
Esta máscara puede tomar cualquier valor entre $0$ (todos los bits son Device ID) y $32$ (todos los bits son Network ID).


Las máscaras de red se presentan a las personas de dos formas principales:\\[-0.5cm]
\begin{itemize}
  \item \otherBase{/N}, donde $N$ es el número de bits asignados al Network ID. 
    Se usa en el formato \concept{CIDR}.
  
  \item El formato \conceptRef{quad decimal}{quad decimal} (\otherBase{A.B.C.D}) de 
    $N$ bits puestos a \otherBase{1}, seguidos de $(32-N)$ bits puestos a \otherBase{0}.
\end{itemize}

\begin{exercise}
La máscara de red del ejemplo anterior puede expresarse en binario 
(aunque no es muy cómodo):
\begin{center}
\begin{bytefield}[bitwidth=1em]{32}
\bitheader{0-31} \\
\bitboxes{1}{1111111111} &
\bitboxes{1}{0000000000 0000000000 00} \\
\bitbox{10}{$10$~bits} & \bitbox{22}{$32-N = 22$~bits}
\end{bytefield}
\end{center}

\begin{itemize}
\item Expresa la máscara de red del ejemplo anterior en los formatos CIDR y cuádruple decimal.
\item ¿Cuál es tu IP y tu máscara de red? Exprésalas de esas mismas dos formas principales.
\end{itemize}
\end{exercise}

\begin{remark}
Las máscaras de red también se usan en IPv6, pero con un valor máximo de $N$ de $128$ en lugar de $32$.
\end{remark}

\subsection{Bloques de direcciones IP}\label{sec:layer3:ip_blocks}

Las máscaras de red se usan combinadas con direcciones IP para definir 
\conceptRef{IP address block}{bloques de direcciones IP} 
que empiezan por los mismos $N$ bits (donde $N$ es la máscara de red).
% 
Por ejemplo, \otherBase{192.168.3.0/24} (en formato \concept{CIDR}) 
y \otherBase{192.168.3.0/255.255.255.0} se refieren al bloque 
de direcciones entre \otherBase{192.168.3.0} y \otherBase{192.168.3.255},
ambas incluidas
(\otherBase{192.168.3.*} en la notación informal de la tabla anterior).

\begin{exercise} 
Los bloques IP pueden ser menores, iguales o mayores que una red IP.
Da ejemplos en los que un bloque IP se refiera a:
\begin{itemize}
\item Una única IP.
\item La mitad de las direcciones IP de la red \otherBase{192.168.3.0/24}.
\item Todas las direcciones de la red \otherBase{192.168.3.0/24}, más 
      $256$ direcciones adicionales.
\item Todas las IP posibles.
\end{itemize}
\end{exercise}


\begin{exercise}\ \\[-0.65cm]
\begin{itemize}
\item ¿Puede expresarse toda \conceptRef{netmask}{máscara de red} como una dirección IP?
\item Enumera todas las máscaras de red posibles.\\[-0.5cm]
\end{itemize}
\end{exercise}

Cuando nos referimos a un bloque de direcciones (como una LAN), 
la dirección IP de esta combinación IP + máscara de red
debe tener todos los bits de Device ID puestos a \otherBase{0}. 
Así, \otherBase{192.168.3.0/24} se refiere correctamente a una red, pero
\otherBase{192.168.3.1/24} se refiere en cambio a un dispositivo 
(el que tiene la IP \otherBase{192.168.3.1}) de esa red (\otherBase{192.168.3.0/24}).

\subsection{LANs y máscaras de red}\label{sec:layer3:lans_and_netmasks}

Las máscaras de red \textit{no} tienen por qué ser múltiplos de $8$. Considera el siguiente ejemplo en el que \\
\otherBase{192.168.76.1/19} (la IP + máscara de red de un dispositivo) se usa para obtener su dirección de red\\
(\otherBase{192.168.64.0/19}):\\

\begin{center}
\begin{bytefield}[bitwidth=1.1em]{32}
\bitheader{0-31}\\
\begin{rightwordgroup}{Device's IP}
\bitbox{8}{192} & \bitbox{8}{168} & \bitbox{8}{76} &\bitbox{8}{1} \\
\bitboxes{1}{11000000} & \bitboxes{1}{10101000} & \bitboxes{1}{0100{\bf1}{\bf1}00} & \bitboxes{1}{0000000{\bf1}}
\end{rightwordgroup} \\

\bitboxes[]{1}{\&\&\&\& \&\&\&\& \&\&\&\& \&\&\&\& \&\&\&
               {\color{color1}\&} 
               {\color{color1}\&}{\color{color1}\&}{\color{color1}\&}{\color{color1}\&}
               {\color{color1}\&}{\color{color1}\&}{\color{color1}\&}{\color{color1}\&}
               {\color{color1}\&}{\color{color1}\&}{\color{color1}\&}{\color{color1}\&}} \\

\begin{rightwordgroup}{Netmask}
\bitboxes{1}{1111111111111111111
   {\color{color1}0}{\color{color1}0}{\color{color1}0}{\color{color1}0}
   {\color{color1}0}{\color{color1}0}{\color{color1}0}{\color{color1}0}
   {\color{color1}0}{\color{color1}0}{\color{color1}0}{\color{color1}0}
   {\color{color1}0}}
\end{rightwordgroup} \\

\bitboxes[]{1}{==== ==== ==== ==== ===
               {\color{color1}=} 
               {\color{color1}=}{\color{color1}=}{\color{color1}=}{\color{color1}=}
               {\color{color1}=}{\color{color1}=}{\color{color1}=}{\color{color1}=}
               {\color{color1}=}{\color{color1}=}{\color{color1}=}{\color{color1}=}} \\

\begin{rightwordgroup}{Network IP}
\bitboxes{1}{11000000} & \bitboxes{1}{10101000} & \bitboxes{1}{0100{\bf0}{\bf0}00} & \bitboxes{1}{0000000{\bf0}} \\
\bitbox{8}{192} & \bitbox{8}{168} & \bitbox{8}{64} &\bitbox{8}{0} 
\end{rightwordgroup}\\
\end{bytefield}
\end{center}

\begin{exercise}\ \\[-0.5cm]
\begin{itemize}
\item En el ejemplo anterior, en realidad solo necesitábamos transformar un valor decimal a binario,
y un valor decimal a binario. ¿Cuál?
\item ¿Es esto siempre cierto al convertir combinaciones IP+máscara de red en direcciones de red?
\item Expresa la máscara de red en formato \conceptRef{quad decimal}{quad decimal}.
\end{itemize}
\end{exercise}



Dado un bloque de direcciones IP + máscara de red y la dirección IP de un dispositivo, es inmediato determinar 
si la IP de ese dispositivo pertenece al bloque (si y solo si coinciden los primeros $N$ bits,
donde \otherBase{/N} es la máscara de red en formato CIDR).

\begin{exercise}
Determina si cada una de las siguientes direcciones pertenece a la misma red que
el siguiente dispositivo \otherBase{123.45.67.89/14}:
\begin{itemize}
\item \otherBase{23.45.67.89}
\item \otherBase{123.45.203.25}
\item \otherBase{123.43.255.255}
\item \otherBase{0.45.67.1}
\end{itemize}
\end{exercise}

Una red IP (una LAN) es un bloque de direcciones, pero no todas esas direcciones 
pueden asignarse a dispositivos. Siempre hay dos direcciones reservadas:
\begin{itemize}
\item La dirección con todos los bits de nodo puestos a \otherBase{0}: identifica la propia LAN.
\item La dirección con todos los bits de nodo puestos a \otherBase{1}: identifica el broadcast (difusión) dentro de la LAN.
\end{itemize}

\begin{exercise}\ \\[-0.5cm]
\begin{itemize}
\item Enumera todas las direcciones que pueden asignarse a dispositivos en \otherBase{192.168.54.0/29}.
\item Indica la dirección de broadcast de esa red.
\end{itemize}
\end{exercise}

La máscara de red \otherBase{/N} determina el número de dispositivos que pueden estar conectados al mismo tiempo
a una red: $2^{32-N} - 2$. Cuando se dice que una LAN es ``pequeña'' o ``más grande'', el tamaño se refiere a ese 
número de direcciones de dispositivo disponibles. 

A la inversa, el número $M$ de dispositivos que necesitamos conectar a la misma LAN determina el 
tamaño de LAN más pequeño (la máscara de red \otherBase{/N} más larga) en el que caben todos esos dispositivos a la vez:

\begin{center}
\begin{tabular}{ccc}
\toprule
\textbf{Máscara de red} & \textbf{Direcciones totales} & \textbf{Dispositivos máx.} \\
\otherBase{/N} & $2^{32-N}$ & $2^{32-N} - 2$ \\
\toprule
\otherBase{/30} & 4 & 2 \\
\otherBase{/29} & 8 & 6 \\
\otherBase{/28} & 16 & 14 \\
\otherBase{/27} & 32 & 30 \\
\otherBase{/26} & 64 & 62 \\
\otherBase{/25} & 128 & 126 \\
$\vdots$ & $\vdots$ & $\vdots$ \\
\bottomrule
\end{tabular}
\end{center}

\begin{remark}
Todos los routers (encaminadores) necesitan una dirección IP válida por cada LAN a la que están conectados.
Por tanto, los routers \textit{sí} cuentan para el número máximo de dispositivos conectados en una LAN.
\end{remark}

\begin{exercise} \ \\[-0.5cm]
\begin{itemize}
\item ¿Cuál es el tamaño mínimo de red en el que caben $256$ dispositivos conectados?
\item ¿Cuál es la máscara de red \otherBase{/N} más larga que permitiría conectarse a todo el mundo en tu ciudad?
\item ¿Cuántas direcciones contiene una red \otherBase{/0}? 
\item ¿Cuántas redes \otherBase{/0} pueden existir? ¿Reciben algún nombre?
\end{itemize}
\end{exercise}


\begin{exercise}
El siguiente fragmento de código toma una dirección IP y muestra sus $4$~bytes expresados en 
cuádruple decimal y en binario.
Amplía (o sustituye) el código para que puedas:
\begin{itemize}
\item Transformar $32$ bits en una cadena en cuádruple decimal.
\item Mostrar una máscara de red en formato de dirección IP y en formato CIDR.
\item Dada la IP y la máscara de red de un dispositivo, determinar si otra IP pertenece a la misma red.
\end{itemize}
\begin{center}
\showCode{snippets/ipcalculator.py}
\end{center}
\end{exercise}




% \todo{exercise netmask calculator}
% 
% \todo{HERE HERE HERE: add and point to the subnetting area}







\section{IP -- Internet Protocol}\label{sec:layer3:ip}

\subsection{Formato del paquete}\label{sec:layer3:ip:format}

\concept{IP} define un formato de paquete con una cabecera 
de \textit{al menos} $20$~bytes. Pueden incluirse partes opcionales,
siempre que el tamaño total de la cabecera sea múltiplo de $32$~bits ($4$~bytes).
El payload puede estar vacío, así que el tamaño mínimo de un datagrama es de $20$~bytes.\\

\begin{center}
\small
\begin{bytefield}[bitheight=3.4em,bitwidth=1.1em]{32}
\bitheader{0-31}\\
\begin{rightwordgroup}[curlyshrinkage=10pt]{Mandatory Header\\($20$~bytes)}
  \bitbox{4}{IP\\Version} 
  & \bitbox{4}{IHL\\{\scriptsize(HLen/$4$)}} 
  & \bitbox{8}{Type of Service}
  & \bitbox{16}{Datagram Length\\{\scriptsize(Head + Payload, bytes)}} \\
  % 
  \bitbox{16}{Fragment ID} & \bitbox{1}{\otherBase{0}}
  & \bitbox{1}{\rotatebox{90}{\scriptsize DF flag}}
  & \bitbox{1}{\rotatebox{90}{\scriptsize MF flag}}
  & \bitbox{13}{Fragment Offset bytes/$8$} \\
  % 
  \bitbox{8}{Time to Live\\(TTL)}
  & \bitbox{8}{Payload protocol}
  & \bitbox{16}{IP Header Checksum}\\
  % 
  \bitbox{32}{Source IP Address} \\
  \bitbox{32}{Destination IP Address}
\end{rightwordgroup} \\
%   
\begin{rightwordgroup}[curlyshrinkage=10pt]{Optional Header\\{\scriptsize($4$ byte words)}}
\bitbox[tlr]{32}{Options\\{\scriptsize(variable length, optional)}} \\
\bitbox[lbr]{24}{} &
\bitbox{8}{Padding to\\$32$ bits}
\end{rightwordgroup} \\
% 
\begin{rightwordgroup}[curlyshrinkage=10pt]{Payload}
\bitbox[tlr]{32}[bitheight=3em]{Payload data\\{\scriptsize(variable length, byte aligned)}} \\
\bitbox[lbr]{24}{\hspace{3cm}\raisebox{0.25cm}{$\vdots$}} & \bitbox[lt]{8}{}
\end{rightwordgroup}
\end{bytefield}
\end{center}

\begin{exercise} \ \\[-0.5cm]
\begin{itemize}
\item ¿Cómo puede calcularse la longitud del payload a partir de la cabecera?
\item ¿Cuál es la longitud máxima del payload?
\item ¿Qué valores pueden tomar los flags de la cabecera IP?
\item ¿Qué es más largo, una dirección IP o una dirección MAC de Ethernet?
\item ¿Es posible enviar exactamente 17 bits de datos de payload en un datagrama IP?
\item ¿Cómo podemos saber si un datagrama encapsula datos TCP o UDP?
\item ¿Por qué la parte opcional de la cabecera debe ser múltiplo de $32$~bits?
\end{itemize}
\end{exercise}

\subsection{Funcionamiento}\label{sec:layer3:ip:operation}

Cuando un dispositivo A quiere enviar un datagrama IP a un dispositivo B, las IP de origen y 
destino de la cabecera son fijas y no cambian para ese datagrama.

Si los dispositivos A y B pertenecen a la misma LAN, el datagrama se envía directamente
(\eg\ encapsulado en un frame (trama) Ethernet). Sin embargo, IP puede llegar mucho más lejos.
Considera el siguiente escenario, en el que A quiere comunicarse con E.

\begin{center}
 \includegraphics[width=0.5\linewidth]{ip_operation_example.pdf}
\end{center}


Aquí, los dispositivos de origen y destino pertenecen a LANs distintas.
El datagrama se pasa primero al router del origen.
Después, cada router intermedio lo pasa al siguiente hasta que el datagrama llega 
al último router, que está conectado a la LAN de destino. Por último, ese último router 
entrega el datagrama al destino. 

\begin{center}
\includegraphics{ip_operation.eps}
\end{center}

Las direcciones IP de origen y destino no cambian a lo largo de todo el trayecto.
En cambio, las direcciones MAC cambian en cada \conceptRef{hop}{salto} de LAN. 

\section{Encaminamiento IP}\label{sec:layer3:routing}

Los routers son los dispositivos encargados de reenviar los datagramas IP
mirando la dirección IP de destino de cada uno, presente en la cabecera IP. Los routers no necesitan 
desencapsular las cabeceras de la capa 4 y superiores. Por tanto, sus pilas de red solo necesitan llegar hasta la capa~3
(y por eso a veces se les llama dispositivos de capa~3).

\begin{center}
\includegraphics[width=0.9\linewidth]{forward_layers.pdf}
\end{center}


Cuando un router detecta un frame de red $F$ en una de sus \conceptRef{interface}{interfaces} (tarjetas de red),
intenta reconocer las siguientes partes de un datagrama.
\begin{center}
\begin{bytefield}{48}
\bitbox{12}{Layer~2 header} & \bitbox{12}{IP header} & \bitbox{24}{Datagram payload}
\end{bytefield}
\end{center}
Si lo reconoce, el router procesa el datagrama de la siguiente manera:
\begin{enumerate}
\item Se descarta el frame $F$ completo salvo que la MAC de destino (en la cabecera de capa~2) 
  sea la de la interfaz del router que ha detectado el frame.\\[-0.25cm]
  
\item El router mira la dirección IP de destino (en la cabecera IP de $F$) y su \conceptRef{routing table}{tabla de enrutamiento} 
  para decidir el siguiente \conceptRef{hop}{salto} del datagrama, \ie, el dispositivo al que este router le pasará 
  el datagrama, y qué \conceptRef{interface}{interfaz} usar (véase la Sección~\ref{sec:layer3:routing} más abajo).\\[-0.25cm]
  
\item El router crea y envía un frame $F'$ haciendo los dos cambios siguientes sobre $F$:\\[-0.25cm]
  \begin{enumerate}[label=\alph*)]
  \item La cabecera de capa~2 de $F'$ cambia por completo:
    \begin{itemize} 
    \item La MAC de destino es la del siguiente salto,
    y la MAC de origen es la del router que reenvía el datagrama.
    La IP del siguiente salto la da la \conceptRef{routing table}{tabla de enrutamiento} (véase la Sección~\ref{sec:layer3:routing}), 
    y su MAC se obtiene a partir de esa IP mediante ARP (Sección~\ref{sec:layer3:arp}).\\[-0.4cm]
    %     
    \item La MAC de origen del frame $F'$
    es distinta de la MAC de destino del frame $F$ recibido si el router 
    recibió $F$ por una interfaz (\eg\ \inlineCode{eth7}) y lo reenvía por otra 
    (\eg\  \inlineCode{eth9}).
    \end{itemize}
    
    \begin{center}     
    \vspace{0.15cm}
    \includegraphics[width=0.8\linewidth]{router_ports.jpg}
    \vspace{0.15cm}
    \end{center}
    
  \item La cabecera IP de $F'$ difiere de la de $F$ en lo siguiente:
    \begin{itemize}
    \item El campo \concept{TTL} (tiempo de vida) se decrementa en $1$.\\
    Si el TTL llega a $0$, el datagrama se descarta en lugar de reenviarse.
    \item Cuando el campo Options está presente, a menudo se actualiza.
    \item El \concept{checksum} (suma de verificación) de la cabecera se actualiza en consecuencia.\\
    \end{itemize}
  \end{enumerate}
\end{enumerate}

\begin{remark}
 Fíjate en cómo las direcciones IP y el payload se copian sin modificar
  cuando se reenvía un datagrama.
\end{remark}




Cuando un router recibe un datagram IP para reenviarlo y ese router 
está conectado a la misma LAN que la dirección IP de destino, el router realiza una \conceptRef{direct/indirect delivery}{entrega directa} (DD).
% 
En caso contrario, el datagrama se reenvía al siguiente router en una \conceptRef{direct/indirect delivery}{entrega indirecta} (ID).
% 
Un router sabe si debe realizar DD o ID, así como qué dirección IP de destino (y de origen) usar,
gracias a su \conceptRef{routing table}{tabla de enrutamiento}.

Todas las tablas de enrutamiento tienen al menos las siguientes columnas:
\begin{itemize}
\item \textbf{Bloque de direcciones de destino}: una o más direcciones IP que pueden expresarse como una combinación IP + máscara de red.
  Una fila de la tabla solo puede aplicarse a un datagrama si su IP de destino está dentro del bloque de direcciones de destino de esa fila:
  \begin{itemize}
  \item Si solo se aplica una fila de la tabla, esa fila se usa para el encaminamiento.
  \item Si se aplica más de una fila a un datagrama, se usa la fila con el bloque de direcciones de destino más pequeño 
    (la más específica).
  \end{itemize}
  Un bloque de destino \conceptRef{default routing table}{\textit{default}} indica el bloque \otherBase{0.0.0.0/0}.
  Como este es el bloque más grande posible, solo se aplica si no se aplica ninguna de las otras filas. 
  La entrada \textit{default} no es obligatoria.\\[-0.3cm]
  
\item \textbf{Gateway/via} (pasarela): si la fila corresponde a una entrega indirecta, esta columna contiene la dirección IP
  del siguiente router. En caso contrario, esta columna contiene \otherBase{0.0.0.0} para indicar una entrega directa.\\[-0.3cm]
  
\item \textbf{Interfaz}: nombre interno de la interfaz de red (\eg\ de la tarjeta Ethernet) que se usará para entregar el 
  datagrama. Los sistemas operativos pueden usar la convención que quieran, siempre que puedan identificarse de forma única en la tabla.
  Los nombres habituales empiezan por \inlineCode{eth}, \inlineCode{wlan}, \inlineCode{atm}, 
  y más recientemente \inlineCode{en} y \inlineCode{wl}.
\end{itemize}

\begin{remark}
Todos los dispositivos tienen su propia tabla de enrutamiento, no solo los routers.
\end{remark}

\begin{exercise}
A continuación se presenta la \conceptRef{routing table}{tabla de enrutamiento} de un router pequeño.

\begin{center}
\vspace{0.2cm}
\begin{tabular}{ccc}
\toprule
\textbf{Bloque de direcciones de destino} & \textbf{gateway/via} & \textbf{Interfaz} \\
\toprule
\otherBase{\textit{default}} & \otherBase{192.168.0.1} & \inlineCode{eth0}  \\
\otherBase{192.168.0.0/24} & \otherBase{0.0.0.0} & \inlineCode{eth0} \\
\otherBase{172.16.0.4/32} & \otherBase{0.0.0.0} & \inlineCode{eth1} \\
\otherBase{10.49.4.0/28} & \otherBase{10.49.4.1} & \inlineCode{eth2} \\
\otherBase{10.49.4.3/32} & \otherBase{10.49.4.2} & \inlineCode{eth2} \\
\bottomrule
\end{tabular}
\vspace{0.2cm}
\end{center}

\begin{itemize}
\item Dibuja un mapa de red coherente con la tabla de enrutamiento, que contenga este router.
  Para las interfaces del router, indica su nombre y unas direcciones IP y de capa~2 válidas.\\[-0.3cm]
  
\item ¿Cómo reenviaría este router un datagrama dirigido a una dirección IP pública?\\[-0.3cm]

\item Da una dirección IP de destino que se procesaría usando la fila \otherBase{10.49.4.0/28},
y otra que se procesaría usando la fila \otherBase{10.49.4.2/32}. 
Piensa en un caso de uso para tener estas dos reglas.
\end{itemize}
\end{exercise}

Para que Internet funcione, los routers que interconectan LANs a escala global deben estar bien configurados. 
A veces no es así, lo que puede provocar errores de encaminamiento. Dos problemas importantes 
que uno puede encontrarse son:
\begin{itemize}
\item ``Network is unreachable'': ninguna de las filas de la tabla de enrutamiento usa un bloque de direcciones que contenga la 
  dirección de destino del datagrama. 
  El router no puede saber cuál debería ser el siguiente salto, así que el datagrama no se reenvía.
  Esto provoca un error ``ICMP Destination Unreachable'' (véase la Sección~\ref{sec:layer3:icmp}).
  
\item ``Time to Live exceeded'': cada vez que se reenvía un datagrama,
  su TTL se decrementa, como se mencionó en la Sección~\ref{sec:layer3:ip:format}.
  Si varios routers están configurados de forma que el datagrama siga algún tipo de ciclo (bucle), 
  el datagrama se reenviaría infinitamente.
  Cuando el TTL llega a 0, el datagrama se descarta y se produce este error.
  Esto provoca un error ``ICMP Time Exceeded'' (véase la Sección~\ref{sec:layer3:icmp}).
\end{itemize}

\begin{exercise}
¿Puede producirse el error ``Network unreachable'' si un router tiene una entrada \otherBase{\textit{default}}?
\end{exercise}

\begin{exercise}
Propón un ejemplo de LANs interconectadas y especifica las tablas de enrutamiento necesarias para que un datagrama 
siga un ciclo y produzca el error ``TTL exceeded''.
\end{exercise}

\section{Subnetting IP}\label{sec:layer3:subnetting}
\begin{bytefield}[bitwidth=1.35em]{32}
\bytefieldsetup{bitheight=8.4em}
\bitbox[]{1}{\raisebox{2.25em}{\rotatebox{90}{$\leftarrow$\otherBase{0.0.0.0}}}} &
\bitbox[]{7}{} &
\bitbox[]{1}{\raisebox{2.5em}{\rotatebox{90}{$\leftarrow$\otherBase{64.0.0.0}}}} &
\bitbox[]{7}{} &
\bitbox[]{1}{\raisebox{3.25em}{\rotatebox{90}{$\leftarrow$\otherBase{128.0.0.0}}}} &
\bitbox[]{7}{} &
\bitbox[]{1}{\raisebox{3em}{\rotatebox{90}{$\leftarrow$\otherBase{192.0.0.0}}}}  &
\bitbox[]{6}{} &
\bitbox[]{1}{\raisebox{7em}{\rotatebox{90}{$\leftarrow$\otherBase{255.255.255.255}}}}  \\
\bytefieldsetup{bitheight=4ex}
\bitbox{32}{\otherBase{0.0.0.0/0}} \\
\bitbox{16}{\otherBase{0.0.0.0/1}} & \bitbox{16}{\otherBase{128.0.0.0/1}} \\
\bitbox{8}{\otherBase{0.0.0.0/2}} & \bitbox{8}{\otherBase{64.0.0.0/2}} & \bitbox{8}{\otherBase{128.0.0.0/2}} & \bitbox{8}{\otherBase{192.0.0.0/2}} \\
% 
\bitbox{4}{\otherBase{/3}} & \bitbox{4}{\otherBase{/3}} & \bitbox{4}{\otherBase{/3}} & \bitbox{4}{\otherBase{/3}} &
\bitbox{4}{\otherBase{/3}} & \bitbox{4}{\otherBase{/3}} & \bitbox{4}{\otherBase{/3}} & \bitbox{4}{\otherBase{/3}} \\

\bitbox{2}{\otherBase{/4}} & \bitbox{2}{\otherBase{/4}} & \bitbox{2}{\otherBase{/4}} & \bitbox{2}{\otherBase{/4}} & 
\bitbox{2}{\otherBase{/4}} & \bitbox{2}{\otherBase{/4}} & \bitbox{2}{\otherBase{/4}} & \bitbox{2}{\otherBase{/4}} &
\bitbox{2}{\otherBase{/4}} & \bitbox{2}{\otherBase{/4}} & \bitbox{2}{\otherBase{/4}} & \bitbox{2}{\otherBase{/4}} & 
\bitbox{2}{\otherBase{/4}} & \bitbox{2}{\otherBase{/4}} & \bitbox{2}{\otherBase{/4}} & \bitbox{2}{\otherBase{/4}} \\ 
% 
\bitbox{1}{\raisebox{0.25em}{$\vdots$}} & \bitbox{1}{\raisebox{0.25em}{$\vdots$}} & \bitbox{1}{\raisebox{0.25em}{$\vdots$}} & \bitbox{1}{\raisebox{0.25em}{$\vdots$}} & 
\bitbox{1}{\raisebox{0.25em}{$\vdots$}} & \bitbox{1}{\raisebox{0.25em}{$\vdots$}} & \bitbox{1}{\raisebox{0.25em}{$\vdots$}} & \bitbox{1}{\raisebox{0.25em}{$\vdots$}} & 
\bitbox{1}{\raisebox{0.25em}{$\vdots$}} & \bitbox{1}{\raisebox{0.25em}{$\vdots$}} & \bitbox{1}{\raisebox{0.25em}{$\vdots$}} & \bitbox{1}{\raisebox{0.25em}{$\vdots$}} & 
\bitbox{1}{\raisebox{0.25em}{$\vdots$}} & \bitbox{1}{\raisebox{0.25em}{$\vdots$}} & \bitbox{1}{\raisebox{0.25em}{$\vdots$}} & \bitbox{1}{\raisebox{0.25em}{$\vdots$}} & 
\bitbox{1}{\raisebox{0.25em}{$\vdots$}} & \bitbox{1}{\raisebox{0.25em}{$\vdots$}} & \bitbox{1}{\raisebox{0.25em}{$\vdots$}} & \bitbox{1}{\raisebox{0.25em}{$\vdots$}} & 
\bitbox{1}{\raisebox{0.25em}{$\vdots$}} & \bitbox{1}{\raisebox{0.25em}{$\vdots$}} & \bitbox{1}{\raisebox{0.25em}{$\vdots$}} & \bitbox{1}{\raisebox{0.25em}{$\vdots$}} & 
\bitbox{1}{\raisebox{0.25em}{$\vdots$}} & \bitbox{1}{\raisebox{0.25em}{$\vdots$}} & \bitbox{1}{\raisebox{0.25em}{$\vdots$}} & \bitbox{1}{\raisebox{0.25em}{$\vdots$}} & 
\bitbox{1}{\raisebox{0.25em}{$\vdots$}} & \bitbox{1}{\raisebox{0.25em}{$\vdots$}} & \bitbox{1}{\raisebox{0.25em}{$\vdots$}} & \bitbox{1}{\raisebox{0.25em}{$\vdots$}} \\
\end{bytefield}

Las LANs pueden considerarse bloques de direcciones (\eg\ expresados en formato \concept{CIDR})
con dos direcciones reservadas. A las redes accesibles directamente a través de Internet se les debe asignar 
un bloque válido lo bastante grande para conectar todos los dispositivos necesarios. Para asegurarse de que no haya dos dispositivos en Internet
con la misma IP pública, no puede haber dos LANs que contengan la misma dirección
(a menos que la dirección sea privada).

Las $2^{32}$~direcciones disponibles en IPv4 se dividen en bloques mediante el 
\conceptRef{netmask}{enmascaramiento de red}. Los bloques grandes pueden dividirse por la mitad aumentando 
en $1$ la longitud de la máscara de red, como se ve en la figura anterior.
% 
Por ejemplo, el bloque inicial \otherBase{0.0.0.0/0} (todas las direcciones) 
puede dividirse en dos bloques \otherBase{/1}, cada uno con la mitad de las direcciones IPv4 
($2^{31}$ -- todavía demasiado grande para cualquier LAN real).
% 
Este proceso continúa de forma independiente para cada uno de los bloques resultantes, y se detiene cuando el tamaño 
de LAN resultante es el deseado. 

Dado un bloque de direcciones \otherBase{/N} (\eg\ el de una única LAN existente), puede que queramos dividirlo 
en subbloques y tener en su lugar varias LANs más pequeñas (\conceptRef{subnet}{subredes}). 
% 
En las subredes solo podemos usar direcciones del bloque original, pero se pueden reservar bloques
para un uso futuro:

\begin{exercise}
Eres el administrador de la red \otherBase{192.168.0.0/24} de 
un laboratorio de investigación. Te gustaría dividirla en las tres subredes resaltadas más abajo.

\begin{itemize}
\item Proporciona un diagrama de las tres redes, incluyendo un router conectado a las tres
  mediante tres interfaces distintas.\\[-0.4cm]
  
\item Proporciona la tabla de enrutamiento \textit{mínima} de ese router (estas redes
no están conectadas a Internet).\\[-0.4cm]
  
\item ¿Cuántos dispositivos pueden conectarse en total a estas tres redes, incluidos los routers?\\[-0.4cm]

\item ¿Cuántas direcciones de la red original quedan sin asignar?\\[-0.4cm]

\item ¿Cuál es la LAN más grande que podemos crear con las direcciones sin asignar?\\[-0.75cm]
\end{itemize}
% 
\begin{center}
\begin{bytefield}[bitwidth=1.375em]{32}
\bitbox{32}{\otherBase{192.168.0.0/24}} \\
% 
\bitbox{16}{\otherBase{192.168.0.0/25}} 
& \bitbox{16}{\otherBase{192.168.0.128/25}} \\
% 
\bitbox{8}[bgcolor=color1!40!white]{\otherBase{192.168.0.0/26}} 
& \bitbox{8}{\otherBase{192.168.0.64/26}} 
& \bitbox{8}{\otherBase{192.168.0.128/26}} 
& \bitbox{8}{\otherBase{192.168.0.192/26}} \\
% 
\bitbox{4}{\otherBase{/27}} &
\bitbox{4}{\otherBase{/27}} &
\bitbox{4}{\otherBase{/27}} &
\bitbox{4}{\otherBase{/27}} &
\bitbox{4}{\otherBase{/27}} &
\bitbox{4}[bgcolor=color1!40!white]{\otherBase{/27}} &
\bitbox{4}[bgcolor=color1!40!white]{\otherBase{/27}} &
\bitbox{4}{\otherBase{/27}}
\end{bytefield}
\end{center}
\end{exercise}

El proceso de crear varias LANs más pequeñas a partir de un único bloque de direcciones se denomina \conceptRef{subnet}{subnetting} (división en subredes).
El proceso inverso, combinar varias LANs en una mayor, se denomina \conceptRef{supernet}{supernetting} (agregación de redes). 

Dos LANs \otherBase{/N} pueden combinarse si y solo si sus direcciones de red 
(su IP sin máscara de red) comparten los primeros $N-1$ bits. Este proceso puede continuar hasta llegar 
a una red \otherBase{/0}, siempre que poseamos las direcciones IP de ambas mitades de cada supernet.

\begin{exercise}\ \\[-0.5cm]
\begin{itemize}
\item ¿Qué otra red puedes combinar con \otherBase{192.168.0.0/24}
para obtener una red \otherBase{/23} $N_{/23}$? ¿Cuál es su dirección de red?

\item $N_{/23}$ forma parte de exactamente una red \otherBase{/22}, $N_{22}$. 
  Enumera todas las redes \otherBase{/24} que necesitarías poseer para combinarlas en $N_{22}$.
\end{itemize}
\end{exercise}

El \conceptRef{subnetting}{subnetting} suele basarse en cuántas direcciones IP deben estar 
disponibles en cada subred. 
Ese tamaño determina el número de bits de la máscara de red de la subred, como se vio en la Sección~\ref{sec:layer3:lans_and_netmasks}.
Pueden añadirse restricciones adicionales, \eg\ hacer que las direcciones IP de una subred sean las mayores,
o menores que las de otras subredes. 

\begin{remark}
Considera por ejemplo la red \otherBase{192.168.0.0/16}. 
Queremos dividir esta red en subredes y crear la red más pequeña que pueda 
conectar $100$ dispositivos simultáneamente. 

Como se vio en la Sección~\ref{sec:layer3:netmasks}, esto determina que, 
si debe crearse una subred a partir de una  para albergar al menos 100 nodos,
hay que dedicar al menos $7$~bits de dirección ($2^6 - 2 = 62,\,2^7=126$) 
a los Device ID y los $32-7=25$ restantes 
a identificar la red, \ie, la subred puede ser tan pequeña como una \otherBase{/25}.
\ \\
\begin{center}
\begin{bytefield}[bitheight=2em]{32}
\bitheader{0-31} \\
\bitbox{16}{\otherBase{192.168}}
& \bitbox{9}{Free bits}
& \bitbox{7}{Device ID}
\end{bytefield}
\end{center}

Los $9$ bits restantes (posiciones $16$ a $24$) pueden fijarse libremente si no
tenemos ninguna restricción adicional. Elegir un valor de estos bits significa
elegir una de las $2^9$ redes de tamaño \otherBase{/25} contenidas 
en la red original, \otherBase{192.168.0.0/16}.
Poner esos bits libres a \otherBase{0} da el rango de IP más bajo,
y ponerlos a \otherBase{1}, el más alto.

Cuando hay que crear varias subredes, es esencial verificar 
que ninguna subred está contenida en otra. En caso contrario, la misma dirección IP 
pertenecería a dos LANs, y el encaminamiento sería imposible. Por ejemplo,
si hay que crear tres subredes para 100, 500 y 200 dispositivos a partir de \otherBase{192.168.0.0/16},
con rangos de direcciones en ese orden, una posible división sería:

\begin{center}
\begin{bytefield}{32}
\bitheader{0-31}\\
\bitbox{16}{\otherBase{192.168}} & \bitboxes{1}{0000 0000 0} & \bitbox{7}{Device ID}
& \bitbox[]{2}{$\rightarrow$} & \bitbox[]{8}{\otherBase{192.168.0.0/25}} \\
\bitbox{16}{\otherBase{192.168}} & \bitboxes{1}{0000 001} & \bitbox{9}{Device ID}
& \bitbox[]{2}{$\rightarrow$} & \bitbox[]{8}{\otherBase{192.168.2.0/23}} \\
\bitbox{16}{\otherBase{192.168}} & \bitboxes{1}{0000 0001} & \bitbox{8}{Device ID}
& \bitbox[]{2}{$\rightarrow$} & \bitbox[]{8}{\otherBase{192.168.1.0/24}} \\
\end{bytefield}
\end{center}

Recuerda que la dirección de red de cada subred resultante tiene todos los bits de Device ID puestos a \otherBase{0},
y que lo más habitual es presentar esa dirección en formato \conceptRef{quad decimal}{quad decimal}.
\end{remark}

\begin{exercise}\ \\[-0.5cm]
\begin{itemize}
\item En el ejemplo anterior, ¿por qué los bits libres de la segunda subred son \otherBase{0000001} 
y no \otherBase{0000000}?\\[-0.4cm]

\item Repite el ejercicio asignando las direcciones más altas posibles a la tercera,
  luego las siguientes direcciones más altas posibles a la segunda subred, 
  y las direcciones más bajas posibles a la primera subred.\\[-0.4cm]
  
\item ¿Cuál es la IP de broadcast de cada subred?
\end{itemize}
\end{exercise}


% 
\section{Fragmentación IP}\label{sec:layer3:ip_fragmentation}
\label{sec:layer3:fragmentation}

Antes de que un dispositivo (incluidos los routers y los ordenadores normales)
pueda enviar un datagrama por una LAN, debe verificar que su longitud en bytes 
no es mayor que la MTU de esa LAN. En caso contrario, el datagrama debe 
\conceptRef{fragmentation}{fragmentarse} en dos o más datagramas más pequeños que sí quepan en la MTU de la LAN.

Los fragmentos más pequeños son datagramas completos con cabecera y payload. Por tanto, 
más fragmentos implican mayor sobrecarga. Cada uno de esos fragmentos se envía 
y se reenvía individualmente como cualquier otro datagrama, y puede que haya que refragmentarlo 
si atraviesa otra LAN con una MTU aún menor.

Cuando un datagrama necesita fragmentarse, la pila de red sigue 
los siguientes pasos:
\begin{enumerate}

\item Se comprueba el \conceptRef{DF}{flag DF} (don't fragment) del datagrama original.
    Si DF=\otherBase{1} pero se necesita fragmentación, el datagrama no se fragmenta 
    y se produce un error.\\[-0.2cm]

\item Los datos de payload se dividen en trozos de tamaño como mucho $S$, de forma que $S + H = MTU$ 
  ($H$ es la longitud de la cabecera del datagrama, típicamente $20$~bytes). Solo el último trozo
  puede tener un tamaño menor que $S$. Esto garantiza que ningún fragmento supere la MTU,
  que es precisamente el objetivo de la fragmentación.\\[-0.2cm]
  
\item La cabecera del datagrama original se copia en la de los fragmentos con las siguientes modificaciones:
  \begin{itemize}
  \item El campo Datagram ID se fija al mismo valor en todos los fragmentos de un datagrama. 
    El SO se asegura de elegir un ID que su pila de red no haya usado recientemente.
    El campo Datagram ID se ignora en los datagramas no fragmentados.\\[-0.3cm]
    
  \item El campo \conceptRef{offset}{Offset} (desplazamiento) de cada fragmento indica la posición (índice de byte)
    de su primer byte de datos de payload dentro del payload del datagrama antes de la fragmentación.
    Este campo no se expresa en bytes, sino en palabras de $8$~bytes.
    %     
    Dado un offset en bytes, el campo \conceptRef{offset}{Offset} se calcula dividiendo entre $8$.
    Por tanto, los offsets de los fragmentos solo pueden ser múltiplos de $8$.\\[-0.3cm]
    
  \item El \conceptRef{MF}{flag MF} (more fragments) se pone a \otherBase{1} (activado)
    en todos los fragmentos salvo en el último. Esto debe cumplirse incluso si uno o más
    fragmentos se refragmentan en una LAN posterior. Así, al fragmentar un 
    datagrama $D$ con MF=\otherBase{1}, ninguno de los fragmentos resultantes\\[-0.2cm]
  \end{itemize}
  
\item Los datagramas fragmento se envían igual que se enviaría un datagrama normal.
  Los fragmentos pueden refragmentarse, pero solo los reensambla el host (equipo conectado a la red) de destino.
  
\end{enumerate}


\begin{exercise}
Dados el \conceptRef{offset}{Offset} y el flag MF, ¿cómo podemos diferenciar
un datagrama que no ha sido fragmentado de un fragmento de datagrama?
\end{exercise}

\begin{exercise}
¿Por qué sería problemático poner MF=\otherBase{0} en el último datagrama 
producido al fragmentar un datagrama con MF=\otherBase{1}?
\end{exercise}

\begin{exercise}
Un dispositivo necesita enviar un datagrama con una cabecera IP de $20$~bytes y una longitud de payload de $2170$~bytes.
El dispositivo está conectado a LAN1, y el datagrama debe atravesar 
LAN2 y LAN3 hasta llegar a la LAN de destino, LAN4. Las MTU de estas LANs son las siguientes:
\begin{center}
\begin{tabular}{r cccc}
\toprule
\textbf{LAN} & LAN1 & LAN2 & LAN3 & LAN4 \\
\textbf{MTU} & 1500 & 710  & 4096 & 864  \\
\bottomrule
\end{tabular}
\end{center}

\begin{itemize}
\item Indica el offset y la longitud del payload (ambos en bytes),
 así como los flags MF y DF de los $4$ datagramas
 que llegan a LAN4 como resultado del envío del datagrama
 (suponiendo que no hay pérdidas ni duplicaciones).
  
\item ¿Puede el dispositivo de origen conocer las MTU de las LANs intermedias y de 
destino?
\end{itemize}

\end{exercise}


\begin{exercise}
Los datagramas que contienen fragmentos pueden llegar duplicados y/o desordenados. 
¿Cómo debería tratar la pila de red receptora estos datagramas?
\end{exercise}


\section{ARP -- Address Resolution Protocol}\label{sec:layer3:arp}

El Address Resolution Protocol (\concept{ARP}) lo usa la capa~3 (IP)
para obtener la dirección MAC (física) de destino correspondiente a una dirección IP conocida.

Los paquetes ARP se envían como payload de frames de capa~2 (\eg\ Ethernet). 
Los paquetes ARP \textit{no} son datagramas y \textit{no} contienen cabeceras IP.

\begin{remark}
La IP que consulta la capa~3 es la del router en la entrega indirecta y la del dispositivo de destino
en la entrega indirecta. Recuerda también que la IP de destino del datagrama no cambia durante el reenvío.
\end{remark}

\subsection{Formato del paquete}\label{sec:layer3:arp:packet_format}

\vspace{0.2cm}
\begin{center}
\begin{bytefield}[bitheight=3em,bitwidth=1.1em]{32}
\bitheader{0-31}\\
% 
\bitbox{16}{Layer~2 address type \\ (\otherBase{1}: Ethernet)} 
& \bitbox{16}{Layer~3 address type \\ (\otherBase{0x0800}: IP)} \\ 
% 
\bitbox{8}{Layer~2\\address length} & \bitbox{8}{Layer~3\\address length} 
& \bitbox{16}{ARP operation\\(\otherBase{1}: request, \otherBase{2}: reply)} \\
% 
\bitbox[tl]{16}[bgcolor=color2!50!white]{Sender Layer~2 address} &
\bitbox[trb]{16}[bgcolor=color2!50!white]{} \\
\bitbox[lbr]{16}[bgcolor=color2!50!white]{} 
& \bitbox{16}[bgcolor=color2!50!white]{Sender Layer~3 address ...} \\
\bitbox{16}[bgcolor=color2!50!white]{... Sender Layer~3 address}
& \bitbox[ltr]{16}[bgcolor=color1!50!white]{Receiver Layer~2 address} \\
\bitbox[tlb]{16}[bgcolor=color1!50!white]{} & \bitbox[rb]{16}[bgcolor=color1!50!white]{} \\
\bitbox{32}[bgcolor=color1!50!white]{Receiver Layer~3 address}
\end{bytefield}
\end{center}

\begin{remark}
El protocolo ARP \textit{no} usa cabeceras IP.
\end{remark}
\begin{remark}
``Sender'' y ``receiver'' deben interpretarse como ``... de este mensaje ARP''.
\end{remark}


\begin{exercise}\ \\[-0.5cm]
\begin{itemize}
\item ¿Cuál es el propósito de los $4$ primeros campos de un paquete ARP?
\item ¿Qué efecto tienen sobre los campos restantes?
\end{itemize}
\end{exercise}

\begin{exercise}
¿Cuántas direcciones válidas aparecen en un frame Ethernet que transporta una petición ARP?
¿Y en una respuesta ARP?
\end{exercise}


\subsection{Funcionamiento}\label{sec:layer3:arp:operation}

Cuando un nodo \node{A} necesita traducir una dirección IP a una dirección MAC, envía un frame de broadcast
(MAC de destino \otherBase{ff:ff:ff:ff:ff:ff} que contiene un paquete ARP. 
En ese paquete ARP, la operación se fija a \inlineCode{REQUEST} (\otherBase{1}),
la IP de destino es la que se está traduciendo, y
la MAC de destino se fija a \otherBase{00:00:00:00:00:00}.

Cuando una pila de red \node{B} ve un paquete de petición ARP, 
envía un paquete de respuesta ARP solo si 
su IP configurada es la de la dirección IP de destino del ARP.
Esa respuesta es idéntica a la petición con estos cambios:
\begin{itemize}
\item El código OP del paquete ARP se fija a \inlineCode{REPLY} (\otherBase{2}).
\item ``Origen'' se refiere ahora a \node{B} (en la petición se refería a \node{A}),
  y ``destino'' a \node{A}.
\item Las $4$ direcciones del paquete están rellenas en una respuesta 
  (no pueden ser \otherBase{00:00:00:00:00:00}).
\end{itemize}

\begin{exercise}
Cuando se ejecuta el comando \inlineCode{ip route} en un nodo, contiene:

\begin{verbatim}
default via 192.168.0.1 dev wlp5s0 src 192.168.0.2 metric 600 [...]
\end{verbatim}

A su vez, el comando \inlineCode{ip link} contiene las siguientes líneas:

\begin{verbatim}
2: enp42s0: <NO-CARRIER,BROADCAST,MULTICAST,UP> mtu 1500 state DOWN [...]
    link/ether d8:43:ae:61:ed:a4 brd ff:ff:ff:ff:ff:ff
3: wlp5s0: <BROADCAST,MULTICAST,UP,LOWER_UP> mtu 1500 state UP [...]
    link/ether 5c:e4:3a:18:97:1b brd ff:ff:ff:ff:ff:ff 
\end{verbatim}

\begin{itemize}
\item El nodo necesita realizar la entrega indirecta de un datagrama,
pero no conoce la dirección MAC de su router gateway. Indica el valor de 
todos los campos (incluyendo la cabecera Ethernet y el mensaje ARP) de la petición ARP 
que el nodo enviará para obtener la MAC del router.

\item Describe los valores de los campos de la respuesta ARP (puedes usar cualquier dirección
      válida como MAC del router).
\end{itemize}
\end{exercise}

No es probable que la dirección MAC asociada a una IP cambie muy a menudo.
Cuando se recibe una respuesta ARP válida, el SO actualiza una 
\conceptRef{arp cache}{caché} añadiendo (o actualizando) las traducciones IP $\leftrightarrow$ MAC recibidas.
% 
Si el SO necesita enviar un datagrama a esa IP antes de que expire esa entrada de la caché,
se usa directamente la dirección MAC de la tabla sin necesidad de mensajes ARP adicionales.
% 
El comando \inlineCode{ip neighbor} muestra el contenido de esa tabla, \eg\ 
\begin{verbatim}
192.168.0.1 dev wlp5s0 lladdr cc:2d:14:fc:22:60 REACHABLE
\end{verbatim}

\begin{remark}
ARP puede usarse para escanear qué IP están en uso en una LAN.
En particular, un nodo puede usar ARP para verificar que su dirección IP configurada no la está usando nadie más.
\end{remark}

ARP se propuso en una época en la que la seguridad de las redes todavía no era un problema. Como resultado,
no proporciona ningún mecanismo para verificar que una respuesta ARP dada sea correcta.
% 
Esto puede aprovecharse para \conceptRef{poison}{envenenar} la tabla caché ARP de un dispositivo enviando 
a la víctima respuestas ARP \concept{gratuitous} (no solicitadas) con 
asociaciones IP $\leftrightarrow$ MAC falsas.

Los SO modernos suelen ignorar las respuestas ARP no solicitadas, 
y los switches (conmutadores) modernos filtran los mensajes ARP que no coinciden con las
direcciones IP y MAC reales. Aun así, IPv6 prescinde de ARP en favor de 
un protocolo más complejo, Neighbor Discovery (ND).

\begin{exercise}
¿Qué impacto podría tener envenenar la caché ARP de una víctima, 
cambiando la MAC almacenada del router gateway por la MAC del atacante?
\end{exercise}

\section{ICMP -- Internet Control Message Protocol}\label{sec:layer3:icmp}

El Internet Control Message Protocol (\concept{ICMP}) es un protocolo auxiliar
(técnicamente parte de IP) que se usa para probar y depurar
las comunicaciones de (inter)red. 


\subsection{Formato del paquete}\label{sec:layer3:icmp:packet_format}

Los mensajes ICMP se encapsulan como payload de datagramas IP.
En la cabecera IP, el campo Payload protocol se fija a \otherBase{0x01} (``ICMP'').\\[-0.2cm]

\begin{center}
\begin{bytefield}[bitwidth=1.2em]{32}
\bitheader{0-31}\\
\bitbox{32}[bitheight=3em]{IP Header (see \ref{sec:layer3:ip:format})} \\
\bitbox{8}{ICMP type} & \bitbox{8}{ICMP subtype} & \bitbox{16}{ICMP checksum} \\
\bitbox{32}{Rest of ICMP header (mandatory, depends on type)} \\
\bitbox[tlr]{32}[bitheight=3em]{ICMP data\\{\scriptsize(variable length, byte aligned)}} \\
\bitbox[lbr]{24}{\hspace{3cm}\raisebox{0.25cm}{$\vdots$}} & \bitbox[lt]{8}{} \\
\end{bytefield}
\end{center}

Las combinaciones tipo/subtipo ICMP más comunes que probablemente encuentres hoy en día:
\begin{center}
\begin{tabular}{ccc l}
\toprule
\textbf{Tipo ICMP} & \textbf{Subtipo ICMP} & \textbf{¿Error?} & \textbf{Significado} \\
\toprule
$8$ & $0$       & No & Mensaje ICMP Echo Request (\concept{ping}). \\
$0$ & $0$       & No & Mensaje ICMP Echo Reply (respuesta \concept{ping}). \\
$3$ & $0$--$16$ & Sí & Destination Unreachable. El subtipo explica la causa. \\
$11$ & $0$      & Sí & Time to Live Exceeded: el TTL ha llegado a $0$.  \\
\bottomrule
\end{tabular}
\end{center}

Para estos mensajes ICMP de error, el campo de datos ICMP contiene información sobre el 
datagrama que provocó ese mensaje ICMP de error:\\[-0.2cm]

\begin{center}
\begin{bytefield}[bitwidth=1.2em]{32}
\bitheader{0,7,8,15,16,23,24,31}\\
\bitbox{32}[bitheight=3em]{Problematic datagram's complete IP header} \\
\bitbox{32}[bitheight=2.5em]{Problematic datagram's first $8$~payload bytes} \\
\end{bytefield}
\end{center}

Para los mensajes echo request y echo reply, el campo de datos ICMP es arbitrario dentro de los límites
de los datagramas IP. La fragmentación de datagramas ICMP (\eg\ payloads de ping grandes) es posible,
pero muchos routers ignoran los fragmentos ICMP.

\subsection{Funcionamiento}\label{sec:layer3:icmp:operation}

Los mensajes de error ICMP los envían los routers y los hosts normales,
disparados por otros datagramas:
\begin{itemize}
\item La familia de mensajes ``Destination unreachable'' se envía cuando un datagrama no puede 
entregarse a su destino, \eg\ por una \conceptRef{routing table}{tabla de enrutamiento} mal configurada, 
un puerto cerrado (véase la Sección~\ref{sec:layer4}), o simplemente porque el host
de destino no está encendido.

\item El mensaje ``TTL exceeded'' se envía cuando un router recibe un datagrama para 
entrega indirecta con TTL$=1$ (véase la Sección~\ref{sec:layer3:routing}).
\end{itemize}

Para evitar cantidades infinitas de tráfico, \textit{no} se producen errores ICMP
cuando los datagramas causantes son:
\begin{itemize}
\item Datagramas ICMP (de error o no).
\item Fragmentos de datagrama, salvo el primero.
\end{itemize}

\begin{exercise}
¿Cómo podemos diferenciar el primer fragmento de un datagrama de los fragmentos siguientes?
\end{exercise}

